%!TEX root = ../../rigorous_dyson_theorem.tex

\section{Proof Strategy for Main Theorem}

We now prove Theorem \ref{thm:main} in several steps.

\subsection{Exact resolvent expansion}

From Lemma \ref{lem:resolvent-identity} applied to $\mcL = \mcL_0 - \mcX$, we have the exact identity:
\begin{equation}\label{eq:exact-expansion}
\mcR = \mcR_0 + \mcR_0 \mcX \mcR.
\end{equation}

Taking expectations and using that $\mcR_0$ is deterministic:
\begin{equation}\label{eq:expectation-expansion}
\E[\mcR] = \mcR_0 + \mcR_0 \E[\mcX \mcR].
\end{equation}

\textbf{The key challenge} is to compute $\E[\mcX \mcR]$ in terms of $\E[\mcR]$ alone, closing the equation.

\subsection{Gaussian integration by parts}

We first establish the integration by parts formula for Gaussian random variables.

\begin{lemma}[Stein's Lemma for Matrices]\label{lem:stein}
Let $X_{ij} \sim \mathcal{N}(0, \sigma^2/N)$ be independent Gaussian random variables and let $F : \R^{N \times N} \to \R^{N \times N}$ be differentiable with $\E[\norm{F(X)}_F], \E[\norm{\nabla F(X)}_F] < \infty$. Then
\begin{equation}
\E[X_{ij} F_{kl}(X)] = \frac{\sigma^2}{N} \E\Bigl[\frac{\partial F_{kl}}{\partial X_{ij}}(X)\Bigr].
\end{equation}
\end{lemma}

\emph{(Proof: See Appendix B, Lemma~\ref{lem:stein}).}

\subsection{Self-energy computation}

This is the heart of the proof. We compute $\E[\mcX \mcR]$ using Stein's lemma and isotropy.

\begin{proposition}[Self-Energy Under Isotropy]\label{prop:self-energy}
Let $\mcR = \begin{pmatrix} R_{11} & R_{12} \\ R_{21} & R_{22} \end{pmatrix}$ be the block resolvent. Then
\begin{equation}\label{eq:self-energy-formula}
\E[\mcX \mcR] = \Sigma[\E[\mcR]] \cdot \E[\mcR] + \text{Error},
\end{equation}
where
\begin{equation}
\Sigma\Bigl[\begin{pmatrix} M_{11} & M_{12} \\ M_{21} & M_{22} \end{pmatrix}\Bigr]
= \begin{pmatrix} 0 & \sigma^2 m_{12} I \\ \sigma^2 m_{21} I & 0 \end{pmatrix},
\quad
m_{ij} = \frac{1}{N}\Tr(M_{ij}),
\end{equation}
and $\norm{\text{Error}}_F = O(1/N)$ with probability $1 - O(e^{-cN})$.
\end{proposition}

\emph{(Proof: See Appendix B, Proposition~\ref{prop:self-energy}).}

\begin{remark}[Isotropy Forces Scalar Structure]
The key observation is that isotropy $\E[XBX^\top] = (\sigma^2/N)\Tr(B)I$ forces the self-energy to be proportional to the identity matrix in each block. This reduces the infinite-dimensional problem of tracking full $N \times N$ matrix structures to a finite-dimensional problem of tracking only two scalar order parameters $(m_{12}, m_{21})$. This dimensional reduction is the crucial simplification that makes the analysis tractable.
\end{remark}

\subsection{Fixed point analysis}

We now prove that the Dyson equation \eqref{eq:dyson} has a unique solution.

\begin{proposition}[Scalar Fixed-Point Formulation]\label{prop:fixed-point}
By Lemma \ref{lem:symmetry}, the symmetry of $\mcL(z)$ implies $m_{12} = m_{21}$.
For each $z \in \C^+$ with $\Im(z) \geq \eta > 0$ satisfying $\eta > 2(1 + \norm{A})\sigma$, there exists a unique $x \in \C$ satisfying the scalar fixed-point equation:
\begin{equation}\label{eq:scalar-fp}
x = K(x),
\end{equation}
where
\begin{equation}\label{eq:K-operator}
K(x) := \left(\bTr\left(\bigl(\mcL_0(z) - \Sigma(x)\bigr)^{-1}\right)\right)_{12},
\end{equation}
with $\Sigma(x) = \begin{pmatrix} 0 & \sigma^2 x I \\ \sigma^2 x I & 0 \end{pmatrix}$.

The diagonal order parameters are then determined by:
\begin{align}
m_{11}(z) &= \frac{1}{N}\Tr\Bigl(\bigl(zI - (A + \sigma^2 x I)(A^\top + \sigma^2 x I)\bigr)^{-1}\Bigr), \label{eq:m11-from-x}\\
m_{22}(z) &= \frac{1}{N}\Tr\Bigl(\bigl(I - \tfrac{1}{z}(A^\top + \sigma^2 x I)(A + \sigma^2 x I)\bigr)^{-1}\Bigr). \label{eq:m22-from-x}
\end{align}

Furthermore, $m_{11}(z)$ satisfies $\Im(m_{11}(z)) > 0$ and $x(z)$ depends continuously on $z$. The Dyson equation \eqref{eq:dyson-fp} is solved by the $2N \times 2N$ block matrix $\mcM(z)$ with blocks $M_{ij} = m_{ij}(z) I_N$ where $m_{12} = m_{21} = x$.
\end{proposition}

\emph{(Proof: See Appendix B, Proposition~\ref{prop:fixed-point}).}

\subsection{Convergence proof}

We now prove Part B of Theorem \ref{thm:main}.

% CERTAINTY SCORE: 3/5 - FLAG FOR HUMAN EXPERT REVIEW
% JUSTIFICATION: Self-consistent error equation E_N = K_N · E_N + R_N is stated but not derived rigorously.
% KEY CONCERNS: (1) The matrix K_N and its properties (||K_N|| < 1) are not explicitly constructed.
%               (2) Step 1 needs detailed derivation showing how E_N satisfies this equation.
%               (3) Coupling to concentration inequality in Step 3 should be made more explicit.
\begin{proof}[Proof of Theorem \ref{thm:main}, Part B]
We need to show $\mcG_N(z) = \bTr(\E[\mcR_N(z)]) \to \bTr(\mcM(z))$ almost surely.

\textbf{Step 1: Self-consistent error equation.} From Proposition \ref{prop:self-energy}, we have
\[
\E[\mcR] = \mcR_0 + \mcR_0 \bigl(\Sigma[\E[\mcR]] \E[\mcR] + \text{Error}\bigr).
\]

Rearranging:
\[
\bigl(\mcL_0 - \Sigma[\E[\mcR]]\bigr) \E[\mcR] = I + \mcR_0 \cdot \text{Error}.
\]

Taking the inverse:
\[
\E[\mcR] = \bigl(\mcL_0 - \Sigma[\E[\mcR]]\bigr)^{-1} + (\text{correction from Error}).
\]

Define the error $E_N = \bTr(\E[\mcR_N]) - \bTr(\mcM)$, where $\mcM$ is the fixed point of Proposition \ref{prop:fixed-point}.

By the structure of the equations, $E_N$ satisfies
\[
E_N = K_N \cdot E_N + R_N,
\]
where:
\begin{itemize}
\item $K_N$ is a $2 \times 2$ matrix (deterministic given $A_N$) with $\norm{K_N} \leq C < 1$ for large $N$
\item $R_N$ is the residual term from the Error in Proposition \ref{prop:self-energy}, with $\norm{R_N} = O(1/N)$
\end{itemize}

\textbf{Step 2: Solving for $E_N$.} From the self-consistent equation,
\[
(I - K_N) E_N = R_N.
\]

For large enough $N$, $I - K_N$ is invertible (since $\norm{K_N} < 1$), giving
\[
\norm{E_N} \leq \norm{(I - K_N)^{-1}} \cdot \norm{R_N} \leq \frac{1}{1 - \norm{K_N}} \cdot O(1/N) = O(1/N).
\]

\textbf{Step 3: Almost sure convergence.} The bound $\norm{E_N} = O(1/N)$ holds in expectation. By the concentration inequality (to be proven in Section 3.6), we also have
\[
\mathbb{P}\bigl(\norm{\mcG_N - \bTr(\mcM)} > t\bigr) \leq Ce^{-cNt^2}.
\]

By the Borel-Cantelli lemma, this implies almost sure convergence.
\end{proof}

\subsection{Concentration bounds}

We establish the concentration of $\mcG_N$ around its expectation.

\begin{lemma}[Concentration of Generalized Green's Function]\label{lem:concentration}
For $z \in \C^+$ with $\Im(z) \geq \eta > 0$,
\[
\mathbb{P}\Bigl(\norm{\bTr(\mcR_N(z)) - \E[\bTr(\mcR_N(z))]}_F > t\Bigr) \leq 2 \exp\Bigl(-\frac{cNt^2\eta^4}{C(\sigma, \norm{A})}\Bigr).
\]
\end{lemma}

\emph{(Proof: See Appendix B, Lemma~\ref{lem:concentration}).}

